% !TEX TS-program = xelatex
% !TEX encoding = UTF-8 Unicode
% !Mode:: "TeX:UTF-8"

\documentclass{resume}

\usepackage{enumitem} % 用于自定义列表样式
\usepackage{hyperref} % 用于创建可点击的链接
\hypersetup{colorlinks=true, urlcolor=blue} % 设置链接颜色

% 定义一个新的 project 环境
% #1: 项目标题, #2: 项目时间
% --- 緊湊版 Project 環境定義 ---
\newenvironment{compactproject}[2]{
    \vspace{3pt} % 減小專案間的垂直間距
    \noindent\textbf{\textit{#1}} % 專案標題 (移除 \large)
    \hfill
    \textbf{#2} % 專案日期
    \par\vspace{2pt} % 減小標題與描述間的距離
    % itemsep=0pt 移除了項目間的額外間距
    \begin{itemize}[leftmargin=1.2em, topsep=1pt, partopsep=0pt, itemsep=0pt, parsep=1pt]
}{
    \end{itemize}
}

\usepackage{zh_CN-Adobefonts_external} % Simplified Chinese Support using external fonts (./fonts/zh_CN-Adobe/)
%\usepackage{zh_CN-Adobefonts_internal} % Simplified Chinese Support using system fonts
\usepackage{linespacing_fix} % disable extra space before next section
\usepackage{cite}

\begin{document}
\pagenumbering{gobble} % suppress displaying page number

\name{林俊峰}

% {E-mail}{mobilephone}{homepage}
% be careful of _ in emaill address
\contactInfo{(+86) 18368042405}{linjf21@mails.tsinghua.edu.cn}{AI Infrastructure}{}
% {E-mail}{mobilephone}
% keep the last empty braces!
%\contactInfo{xxx@yuanbin.me}{(+86) 131-221-87xxx}{}
 
% \section{个人总结}
% 本人在校成绩优秀、乐观向上，工作负责、自我驱动力强、热爱尝试新事物，认同开放、连接、共享的Web在未来的不可替代性。在校期间长期从事可视分析(Web的2D/3D时空可视化)相关研究，对Web技术发展趋势及前端工程化解决方案有浓厚兴趣。\textbf{现任职于 BAT 集团。}

% \section{\faGraduationCap\ 教育背景}
\section{教育背景}
\datedsubsection{\textbf{清华大学}, 类脑计算中心, \textit{博士在读}}{2021.9 - 2026.6（预计）}
\ 清华大学综合优秀奖学金
% \ \textbf{排名11/133(前10\%)},中国科学院大学学业奖学金(2次),IEEE Student member,预计2018年6月毕业
\datedsubsection{\textbf{伊利诺伊大学-厄巴纳香槟分校},计算机科学,\textit{访问学者}}{2025.2 - 2025.8}\ 清华大学访问学者奖学金

\datedsubsection{\textbf{清华大学},计算机科学与技术,\textit{工学学士}}{2017.9 - 2021.6}
\ \textbf{GPA 3.66/4.0},国家励志奖学金,学业优秀奖学金,计算机系优秀毕业生


% \section{\faCogs\ IT 技能}
% \section{技术能力}
% % increase linespacing [parsep=0.5ex]
% \begin{itemize}[parsep=0.2ex]
%   \item \textbf{编程语言}: C++, Python
%   \item \textbf{操作系统,数据库与工程构建}: Linux/macOS/MySQL/MongoDB/Git/webpack/Progressive Web App
%   \item \textbf{关键词}: React/Vue.js/D3.js(SVG)/three.js(canvas, WebGL)/chrome extension/Express
% \end{itemize}

% \end{itemize}

\section{学术论文}

\datedsubsection{\textbf{期刊}}{}
\begin{enumerate}
    % \item {HASP: Hierarchical asynchronous parallelism for multi-NN tasks. 
    % Hongyi Li, ..., \textbf{\underline{Junfeng Lin}}, et al. 2023. IEEE Transactions on Computers}. \\(\textbf{TC, CCF-A})

    \item {SongC: A Compiler for Hybrid Near-Memory and In-Memory Many-Core Architecture. 
    \textbf{\underline{Junfeng Lin*}}, Huanyu Qu*, Songchen Ma*, et al. 2023. IEEE Transactions on Computers}. \\ (\textbf{TC, CCF-A, 第一作者})

    \item {A versatile and power-saving brain-inspired computing chip. Yangshu Shen*, Zhenzhi Wu*, Weihao Zhang*, \textbf{\underline{Junfeng Lin*}}, et.al.. 2026.}\\
    (\textbf{Nature}，学生一作，在投)

\end{enumerate}

\datedsubsection{\textbf{会议}}{}
\begin{enumerate}
    \item {CATFlow: Data-Monistic Parallelism Plan Discovery for Distributed Training}. Qu Huanyu*, Zhang Weihao*, \textbf{Junfeng Lin*}, et al.. \\(\textbf{OSDI'26, CCF-A}, 在审)
    \item {VoltanaLLM: Feedback-Driven Frequency Control and State-Space Routing for Energy-Efficient LLM Serving.
    Jiahuan Yu, \textbf{ \underline{Junfeng Lin*}}, et al.}. \\(\textbf{ISCA'26, CCF-A}, 在审)
    
    % \item {MLDSE: Scaling Design Space Exploration Infrastructure for Multi-Level Hardware.
    % Huanyu Qu, Weihao Zhang, \textbf{Junfeng Lin}, et al.}. \\(\textbf{HPCA'26, CCF-A}, 在审)

    % \item {Scalable PIM-based LLM System via Hybrid Complementary PIMs and  In-Transit Computation.
    % Hongyi Li,..., \textbf{Junfeng Lin}, et al.}. \\(\textbf{HPCA'26, CCF-A}, 在审)
    \item {WeiPipe: Weight Pipeline Parallelism for Communication-Effective Long-Context Large Model Training.
    \textbf{\underline{Junfeng Lin*}}, Ziming Liu*, et al.}. \\( \textbf{PPoPP'25, CCF-A, 第一作者})
\end{enumerate}
    



% \begin{onehalfspacing}
% \end{onehalfspacing}

% \datedsubsection{\textbf{DID-ACTE} 荷兰莱顿}{2015年}
% \role{本科毕业设计}{LIACS 交换生}
% 利用结巴分词对中国古文进行分词与词性标注，用已有领域知识训练形成 classifier 并对结果进行调优
% \begin{onehalfspacing}
% \begin{itemize}
%   \item 利用结巴分词对中国古文进行分词与词性标注
%   \item 利用已有领域知识训练形成 classifier, 并用分词结果进行测试反馈
%   \item 尝试不同规则，对 classifier 进行调优
% \end{itemize}
% \end{onehalfspacing}

\section{专利}
% \subsection{\textbf{已授权}}
\begin{enumerate}
        \item 权重流水并行的神经网络分布式训练方法和装置。申请日：2024-07-19。专利号：ZL 2023 1 0160001.1。
    \item 任务处理的方法和装置、众核系统、计算机可读介质。申请日：2023年02月21日。专利号：ZL 2023 1 0160001.1 。
    \item 神经网络任务的代码生成方法和装置。申请日：2022年06月21日。专利号：ZL 2022 1 0705825.8。
\end{enumerate}
% \subsection{\textbf{实质审查}}
% \begin{enumerate}

%     \item 任务编译的方法、编译器、计算机可读介质。申请日：2023年02月21日。公开号：CN119067196A。    \item 时序神经网络的映射方法及装置、加速器。申请日：2022年11月10日。公开号：	
%     CN115660054A。
%     \item 任务映射的方法、电子设备、计算机可读介质。申请日：2023年2月24日。公开号：	
%     CN116382894A。
%     \item 任务处理方法及装置。申请日：2023年1月16日。公开号：	
%     CN116257249A。
% \end{enumerate}

\section{项目经历}
\begin{compactproject}{面向ZeRO的存储通信优化}{2025.9 -- 2025.12}
\item 基于ZeRO3实现进一步ZeRO3.5和ZeRO4策略，在通信受限的场景下相比ZeRO3提升39.1\%的吞吐，相对ZeRO2降低36.6\%的显存
\item 实现Preload-Then-Discard的ZeRO-offload策略，存储受限的场景下提高模型吞吐高达89.4\%

\end{compactproject}
\begin{compactproject}{VotanaLLM: 大模型节能推理系统}{2025.3 -- 2025.6}
\item 设计开发VotanaLLM节能推理框架，通过对Prefill/Decode阶段进行异构频率调度与动态路由，在保证服务质量(SLO)的前提下，将LLM推理能耗最高降低\textbf{36\%}。
\item 相关研究成果已投稿至体系结构顶级会议 \textbf{ISCA'26}。
\end{compactproject}

\begin{compactproject}{权重流水并行框架 (长文本训练优化)}{2023.9 -- 2024.9}
\item 针对低带宽集群的训练瓶颈，提出权重流水(Weight Pipelining)并行方案，通过有效隐藏通信开销显著优化长文本训练性能。
\item 核心成果被并行计算顶级会议 \textbf{PPoPP'25} (CCF-A) 录用。
\end{compactproject}

% \begin{compactproject}{BiMap: 众核类脑芯片编译工具链}{2021.9 -- 2023.9}
% \item 作为核心开发者，为“天机”芯片设计了BiMap编译及仿真工具链。通过构建多级IR(中间表示)架构与编译-仿真协同优化流程，高效解决了混合神经网络的自动化部署难题。
% \item 成果发表于体系结构顶刊 \textbf{IEEE TC} (CCF-A)，项目已在GitHub开源。
% \end{compactproject}

% \begin{compactproject}{强化学习算法研究与应用}{2018.9 -- 2019.12}
% \item 负责强化学习算法研究与实践，基于A2C/SAC实现自动驾驶模型(入选清华“学术推进计划”)。
% \item \textbf{核心成就}：设计的多智能体协作策略，获得NeurIPS 2018 Pommerman全球竞赛\textbf{第七名}。
% \end{compactproject}


\section{实习经历}
\begin{compactproject}
{凯读投资，AI Infra 研究实习生 }{2025.10 -- 2025.12}
\item 探索流水线并行支撑自研分布式训练框架。
\item 研究低精度训练和高效注意力机制在量化模型上的应用。
\end{compactproject}

\begin{compactproject}
{灵汐科技，编译器研发实习生 }{2021.7 -- 2024.9}
\item 作为核心成员，深度参与自研“天机”类脑芯片的全栈工具链研发，包括编译器、仿真器及大模型分布式训练框架的设计与实现。
% \item 将“BiMap编译工具链”及“权重流水并行框架”等核心研究项目成功在“天机”芯片上落地实现，支撑了多篇顶会/顶刊论文的发表。
\end{compactproject}

\begin{compactproject}
{腾讯，软件开发实习生}{2020.7 -- 2020.9}
\item 负责腾讯内部Git平台的WebIDE(在线集成开发环境)的前期技术预研与核心功能开发。
% \item  利用Docker容器化技术，为腾讯Git设计并搭建了一套隔离、安全的在线代码编译与执行环境。
% \item 因表现出色，最终获评“优秀实习项目”。
\end{compactproject}

\section{荣誉奖项}
\textbf{清华大学研究生综合奖学金} \hfill 2025年9月 \\
\textbf{PPoPP 2026 Travel Grant} \hfill 2025年3月 \\
\textbf{清华大学访问学者奖学金} \hfill 2024年11月 \\
\textbf{清华大学研究生综合奖学金} \hfill 2023年9月 \\
\textbf{清华大学计算机系优秀毕业生} \hfill 2021年7月 \\
\textbf{清华大学国家励志奖学金} \hfill 2019年11月 \\
\textbf{清华大学学业优秀奖} \hfill 2019年11月 \\
\textbf{全国部分地区大学生物理竞赛(非物理组)二等奖} \hfill 2018年12月 \\

\end{document}
